\section{One aligned site on two four-leaf trees}
\label{sec:book-phylogeny-toy-tree}

For one alignment column, a phylogenetic model assigns a categorical state to every vertex of a
tree.  The leaves $A,B,C,D$ are observed homologous sequences.  Internal vertices are latent
ancestors whose residue states are summed out.  A biological root, when identified, orients descent
from an ancestor; under a reversible likelihood model, a computational root may simply orient the
recursion.  Each edge carries a substitution transition kernel.  Its branch length usually records
expected substitutions per site under the declared model, which need not equal calendar time or
raw percentage difference.

Two four-leaf topologies expose the role of shared history:
\begin{equation}
\begin{aligned}
\text{star:}&\qquad r\longrightarrow\{A,B,C,D\},\\
\text{binary:}&\qquad
r\longrightarrow\{u,v\},\qquad
u\longrightarrow\{A,B\},\quad v\longrightarrow\{C,D\}.
\end{aligned}
\label{eq:book-star-binary-trees}
\end{equation}
The \emph{patristic distance} $d_{ab}$ is the sum of branch lengths along the unique path between
leaves $a$ and $b$.  Under a stationary reversible single-site process, let $\varphi$ be a centered
eigenfunction whose signal decays at rate $\lambda>0$.  Then
\begin{equation}
\operatorname{Cov}\!\left(\varphi(X_a),\varphi(X_b)\right)
=\operatorname{Var}_{\pi}(\varphi)e^{-\lambda d_{ab}}.
\label{eq:book-tree-mode-covariance}
\end{equation}
This formula is a transparent special case: leaf correlation follows the shared path geometry and
the substitution process together.

For a star with root-to-leaf length $t$, every leaf pair has distance $2t$ and the same covariance
in Eq.~\ref{eq:book-tree-mode-covariance}.  For a balanced binary tree with root-to-internal length
$a$, internal-to-leaf length $b$, and $a+b=t$, the within-clade distances are
$d_{AB}=d_{CD}=2b$, whereas the cross-clade distances are $2t$.  The binary tree therefore creates
stronger correlation within recent clades while retaining weaker correlation across them.  With
unequal branch lengths or another substitution process, the comparison changes; topology alone
does not order all correlations.

Correlated leaves contain less independent information than the same number of independent rows.
Under matched root-to-tip lengths, a star lacks the additional recent clade structure of the binary
example and can yield a larger effective sample size for estimating column frequencies and
covariation.  This improves the visibility of a fixed structural or functional constraint through
reduced phylogenetic redundancy.  It does not increase the underlying residue coupling.

Selection pressure is a property of the genotype-to-demography map, population process, and
substitution rates.  A star or binary topology records the branching history and, together with branch lengths,
the dependence among sampled leaves.  No ranking of selection pressure follows from topology by
itself.  The basic site-independent phylogenetic model applies this construction separately to
each alignment column on a shared tree; residue-interaction models add dependence between columns,
and recombination or gene-tree discordance may make the tree vary along the sequence.
